\scanpage{115}
Let $f(z)=\prod_1^\infty(1+\lambda_i z)$ be an entire function
of \uline{genus} $1$, put $\alpha_i=|\lambda_i|$, $F(z)=\prod_{\unsure{1|i}}^\infty(1+\alpha_i z)$.
Then \uline{$f$ and $F$ have the same order} (which is thus expressed
also as a function of the sequence of the $|\alpha_i|$). Indeed,
the order of $f\leq$ order of $F$ (obvious). On the other
hand, order of $F=$ order of the sequence of the $|\alpha_i|$
(see above) $\leq$ order of $f$ (Hadamard's theorem)
\begin{center}\rule{1.5cm}{0.4pt}\end{center}

One can speak of the order of a \inserted{countable} set of
complex numbers, \deleted{\unsure{independently of}} \inserted{even \unsure{not} arranged in a sequence:}
one takes the infimum of the $s>0$ such that
$\sum|z_i|^s<+\infty$. Rapidly decreasing sequence $=$
of order zero.

Order of a sequence in a Banach space,
or a Fréchet space. (two definitions, according as one
takes strong summability, or absolute summability)

Order of a bounded subset (\deleted{smallest} \inserted{infimum} of the
$s>0$ such that the subset is contained in the\newline
closed convex circled hull of a sequence of order $s$)
\margin{116}
